\documentclass[a4paper]{article}

% Basic packages
% Español (México) con XeLaTeX: fontspec para fuentes Unicode, polyglossia
% para la división silábica y los nombres del documento en la variante mexicana.
\usepackage{fontspec}
\usepackage{polyglossia}
\setdefaultlanguage[variant=mexican]{spanish}
% Los nombres chinos van como «pinyin (original)»: xeCJK da a los caracteres
% Han su propia fuente; sin ella XeLaTeX los omite con un aviso.
% FandolSong no cubre algunos caracteres de nombres (U+73FA); WenQuanYi Zen Hei sí.
\usepackage[AutoFallBack=true]{xeCJK}
\setCJKmainfont{FandolSong-Regular.otf}
\setCJKfallbackfamilyfont{\CJKrmdefault}{WenQuanYi Zen Hei}
% polyglossia no traduce los nombres de listings.
\AtBeginDocument{\providecommand{\lstlistingname}{}\renewcommand{\lstlistingname}{Listado}%
  \providecommand{\lstlistlistingname}{}\renewcommand{\lstlistlistingname}{Índice de listados}}
\usepackage{amsmath, amssymb}
\usepackage{graphicx}
\usepackage[margin=2.5cm]{geometry}
\usepackage[most]{tcolorbox}
\usepackage{etoolbox}
\usepackage{listings}
\usepackage{hyperref}
\usepackage{booktabs}
\usepackage{subcaption}
\usepackage{float}
\usepackage{tikz}
\IfFileExists{pgfplots.sty}{
    \usepackage{pgfplots}
    \pgfplotsset{compat=1.18}
}{}

% Highlight boxes
\newtcolorbox{knowledgebox}[1]{
    enhanced, colback=blue!5!white, colframe=blue!75!black, colbacktitle=blue!75!black,
    coltitle=white, fonttitle=\bfseries, title={#1}, attach boxed title to top left={yshift=-2mm, xshift=2mm},
    boxrule=1pt, sharp corners
}

\newtcolorbox{importantbox}[1]{
    enhanced, colback=yellow!10!white, colframe=yellow!80!black, colbacktitle=yellow!80!black,
    coltitle=black, fonttitle=\bfseries, title={#1}, sharp corners
}

\newtcolorbox{warningbox}[1]{
    enhanced, colback=red!5!white, colframe=red!75!black, colbacktitle=red!75!black,
    coltitle=white, fonttitle=\bfseries, title={#1}, sharp corners
}

% Listings
\lstset{
    language=Python,
    basicstyle=\ttfamily\small,
    keywordstyle=\color{blue},
    stringstyle=\color{red!60!black},
    commentstyle=\color{green!60!black},
    breaklines=true,
    frame=single,
    numbers=left,
    numberstyle=\tiny\color{gray},
    captionpos=b,
    extendedchars=true
}

% Front-page metadata
\newcommand{\notetitle}{MIT 6.S191: Convolutional Neural Networks}
\newcommand{\noteauthors}{Elaboradas a partir de la clase de Alexander Amini}
\newcommand{\notedate}{Primavera de 2025}
\newcommand{\videochannel}{Alexander Amini (MIT)}
\newcommand{\videopublishdate}{2025-03-17}
\newcommand{\videoduration}{61 minutos}
\newcommand{\videourl}{https://www.youtube.com/watch?v=oGpzWAlP5p0}
\newcommand{\videocoverpath}{cover.jpg}
\newcommand{\repourl}{https://github.com/hqhq1025/ai-course-notes}

% Footer with repo info
\usepackage{fancyhdr}
\pagestyle{fancy}
\fancyhf{}
\fancyfoot[C]{\thepage}
\fancyfoot[R]{\footnotesize\color{gray}\href{\repourl}{AI Course Notes}}
\renewcommand{\headrulewidth}{0pt}
\renewcommand{\footrulewidth}{0pt}

\begin{document}

\begin{titlepage}
\centering
{\Large Notas del curso\par}
\vspace{1.2cm}
{\huge\bfseries \notetitle\par}
\vspace{0.8cm}
{\large \noteauthors\par}
\vspace{0.3cm}
{\large \notedate\par}
\vspace{1.2cm}

\ifdefempty{\videocoverpath}{
{\small Al generar las notas, indique la ruta local de la portada del video.\par}
}{
\includegraphics[width=0.82\textwidth,height=0.45\textheight,keepaspectratio]{\videocoverpath}\par
}

\vfill
\begin{tcolorbox}[width=0.9\textwidth, colback=black!2!white, colframe=black!60, sharp corners]
\textbf{Autor o canal del video}: \videochannel\par
\textbf{Fecha de publicación}: \videopublishdate\par
\textbf{Duración del video}: \videoduration\par
\ifdefempty{\videourl}{}{
\textbf{Enlace del video}: \href{\videourl}{\nolinkurl{\videourl}}\par
}\par
\textbf{Página del proyecto}: \href{\repourl}{\nolinkurl{\repourl}}\par
\end{tcolorbox}
\end{titlepage}

\tableofcontents
\newpage

%% --- Inicio del contenido --- %%

\section{Introducción: qué es la visión por computadora}

La visión es uno de los sentidos más importantes del ser humano. Nos permite detectar e interpretar emociones y expresiones faciales, desplazarnos en el mundo, manipular objetos e interactuar con el entorno. Alexander Amini propuso, al abrir esta clase, una definición concisa: la visión por computadora (Computer Vision) es la capacidad de dotar a las computadoras de la habilidad de "saber qué hay y dónde con solo mirar" (\textit{to know what is where by looking}).

\begin{figure}[H]
\centering
\includegraphics[width=0.85\textwidth]{slides-images/slide-02.jpg}
\caption{Definición central de la visión por computadora: "To know what is where by looking." \protect\footnotemark}
\end{figure}
\footnotetext{Diapositiva 2 de las láminas.}

La visión no consiste solo en reconocer objetos de manera estática, sino también en entender los cambios dinámicos de una escena. Por ejemplo, cuando vemos la fotografía de una calle en la ciudad, no solo identificamos automóviles, peatones y semáforos, sino que también inferimos qué automóviles están en movimiento, cuáles están estacionados a un costado y hacia dónde se desplazan los peatones. Esta capacidad de inferir información dinámica a partir de una sola imagen es precisamente lo que la visión por computadora busca lograr.

\begin{figure}[H]
\centering
\includegraphics[width=0.85\textwidth]{slides-images/slide-03.jpg}
\caption{El objetivo de la visión por computadora: descubrir a partir de una imagen qué objetos hay en la escena, dónde están, qué acción está ocurriendo, así como predecir e inferir eventos futuros. \protect\footnotemark}
\end{figure}
\footnotetext{Diapositiva 3 de las láminas.}

\begin{knowledgebox}{¿Por qué es tan difícil la visión por computadora?}
El sistema visual humano ha evolucionado durante millones de años y realiza sin esfuerzo tareas como el reconocimiento de objetos y la comprensión de escenas. Pero para una computadora, incluso un mismo objeto puede verse radicalmente distinto a nivel de píxel debido a factores como el cambio de punto de vista (viewpoint variation), el cambio de escala (scale variation), la deformación (deformation), la oclusión (occlusion), las condiciones de iluminación (illumination conditions), el desorden del fondo (background clutter) y la variación dentro de una misma clase (intra-class variation). Todo esto hace de la visión por computadora un problema sumamente difícil.
\end{knowledgebox}

\subsection{Aplicaciones e impacto de la visión por computadora}

El aprendizaje profundo está impulsando una revolución enorme en el campo de la visión por computadora, y sus aplicaciones están en todas partes:

\begin{figure}[H]
\centering
\includegraphics[width=0.9\textwidth]{slides-images/slide-04.jpg}
\caption{Aplicaciones extensas de la visión por computadora: robótica, asistencia para la accesibilidad, biomedicina, conducción autónoma y computación móvil. \protect\footnotemark}
\end{figure}
\footnotetext{Diapositiva 4 de las láminas.}

\begin{itemize}
    \item \textbf{Detección y reconocimiento facial}: no solo detecta la presencia de un rostro, sino que también analiza puntos clave faciales y microexpresiones. Es una de las primeras tareas de visión por computadora que obtuvo una aplicación extensa.
    \item \textbf{Conducción autónoma}: el laboratorio de Amini en el MIT ha construido sistemas de conducción autónoma de extremo a extremo capaces de conducir de manera autónoma en carreteras nunca antes vistas, usando únicamente entrada visual. A diferencia de la mayoría de los automóviles autónomos, que usan un pipeline modular (varios modelos más mapas predefinidos), el método de extremo a extremo usa un solo modelo que va directamente de la imagen al ángulo del volante.
    \item \textbf{Imagenología médica}: las CNN ya superan la exactitud de radiólogos especializados en la detección temprana de cáncer de mama y son capaces de detectar lesiones que a los médicos humanos se les pasan. Además, se aplican ampliamente en la detección de cáncer de piel, el diagnóstico por imágenes de COVID-19 y otras áreas.
    \item \textbf{Asistencia para la accesibilidad}: por ejemplo, el proyecto Google Project Guideline, que ayuda a personas con discapacidad visual a correr de manera independiente.
    \item \textbf{Computación móvil}: en el teléfono de cada persona se ejecuta una gran cantidad de modelos de visión por computadora.
\end{itemize}

\begin{figure}[H]
\centering
\begin{subfigure}[b]{0.48\textwidth}
\centering
\includegraphics[width=\textwidth]{slides-images/slide-05.jpg}
\caption{Detección facial y reconocimiento de puntos clave}
\end{subfigure}
\hfill
\begin{subfigure}[b]{0.48\textwidth}
\centering
\includegraphics[width=\textwidth]{slides-images/slide-07.jpg}
\caption{Visión por computadora en imagenología médica}
\end{subfigure}
\caption{Aplicaciones representativas de la visión por computadora. \protect\footnotemark}
\end{figure}
\footnotetext{Diapositivas 5 y 7 de las láminas.}

\subsection{Resumen de la sección}

La visión por computadora dota a las computadoras de la capacidad de entender el mundo visual. Aunque para los seres humanos resulta sencillo, para las computadoras es sumamente difícil debido a la enorme variabilidad en la apariencia de los objetos. El auge del aprendizaje profundo está transformando este campo de manera fundamental y ha permitido que las computadoras alcancen, e incluso superen, el nivel humano en muchas tareas visuales.

\newpage
\section{Cómo la computadora «ve» las imágenes}

\subsection{Representación numérica de las imágenes}

Para la computadora, una imagen no es más que números. Una imagen en escala de grises es un arreglo bidimensional en el que cada elemento (píxel) es un entero en el rango $[0, 255]$ que representa el brillo de ese píxel. Una imagen a color, en cambio, es un arreglo tridimensional de tamaño $H \times W \times 3$, donde la tercera dimensión corresponde a los tres canales de color RGB.

\begin{figure}[H]
\centering
\includegraphics[width=0.9\textwidth]{slides-images/slide-12.jpg}
\caption{La imagen es números: una imagen en escala de grises es una matriz bidimensional, y una imagen a color (por ejemplo, de 1080$\times$1080) es un apilamiento de tres matrices de este tipo (los tres canales RGB).\protect\footnotemark}
\end{figure}
\footnotetext{Slides, página 12.}

\begin{importantbox}{Representación matemática de las imágenes}
\begin{itemize}
    \item \textbf{Imagen en escala de grises}: $\mathbf{I} \in \mathbb{R}^{H \times W}$, cada valor de píxel $I_{i,j} \in [0, 255]$
    \item \textbf{Imagen a color}: $\mathbf{I} \in \mathbb{R}^{H \times W \times 3}$, los tres canales corresponden respectivamente a rojo (R), verde (G) y azul (B)
    \item Por ejemplo, una imagen RGB de $1080 \times 1080$ contiene $1080 \times 1080 \times 3 = 3,499,200$ números
\end{itemize}
\end{importantbox}

Amini señala que, desde el punto de vista de la representación, las imágenes en realidad son más fáciles de procesar que el lenguaje, porque las imágenes ya son de manera natural arreglos de números, mientras que el lenguaje requiere primero una conversión, como el embedding de palabras.

\subsection{Dos tipos básicos de tareas en visión por computadora}

Los dos tipos de tareas más básicos en visión por computadora son la \textbf{regresión} (Regression) y la \textbf{clasificación} (Classification):

\begin{figure}[H]
\centering
\includegraphics[width=0.85\textwidth]{slides-images/slide-13.jpg}
\caption{Tareas básicas de la visión por computadora: clasificación (produce una etiqueta de categoría discreta y su probabilidad) y regresión (produce un valor continuo).\protect\footnotemark}
\end{figure}
\footnotetext{Slides, página 13.}

\begin{itemize}
    \item \textbf{Regresión}: la variable de salida toma valores continuos, por ejemplo, predecir el ángulo de giro del volante.
    \item \textbf{Clasificación}: la variable de salida toma etiquetas de categoría discretas, por ejemplo, determinar qué presidente estadounidense es la persona en la imagen. Una tarea de clasificación puede producir una distribución de probabilidad sobre cada categoría.
\end{itemize}

Sin importar de qué tarea se trate, el problema central siempre es la \textbf{detección de características} (Feature Detection) — entender qué características distinguen a las diferentes categorías entre sí.

\subsection{Resumen de la sección}

La computadora representa las imágenes mediante matrices de valores de píxeles. Una imagen a color es un tensor tridimensional (alto $\times$ ancho $\times$ número de canales). Los dos tipos básicos de tareas en visión por computadora —clasificación y regresión— dependen de la capacidad de extraer características de manera efectiva a partir de la imagen.

\newpage
\section{Extracción de características: de diseño manual a aprendizaje automático}

\subsection{Las limitaciones de la extracción manual de características}

Para clasificar imágenes correctamente, necesitamos extraer características que distingan las diferentes categorías. Los métodos tradicionales dependen de la \textbf{extracción manual de características} (Manual Feature Extraction), con el siguiente proceso:

\begin{figure}[H]
\centering
\includegraphics[width=0.9\textwidth]{slides-images/slide-15.jpg}
\caption{Proceso de extracción manual de características: conocimiento del dominio $\rightarrow$ definición de características $\rightarrow$ detección de características para la clasificación. El problema central es que este proceso es extremadamente sensible a los cambios.\protect\footnotemark}
\end{figure}
\footnotetext{Diapositiva 15.}

Por ejemplo, para reconocer un rostro, podrías definir las características como «nariz, ojos, boca». Pero esto plantea de inmediato un problema recursivo: ¿cómo detectar un «ojo»? Necesitarías definir características de nivel aún más bajo para describir la apariencia del ojo.

\begin{warningbox}{El dilema central de las características manuales}
La extracción manual de características enfrenta dos retos fundamentales:
\begin{enumerate}
    \item \textbf{La recursividad en la definición de características}: definir características de alto nivel requiere definir características de nivel más bajo, y la definición de estas últimas requiere características de un nivel aún más bajo, formando una recursión infinita.
    \item \textbf{La variabilidad}: la apariencia de un mismo tipo de objeto varía enormemente según el punto de vista, la iluminación, la escala, la deformación y la oclusión; las reglas hechas a mano no pueden enumerar todos los casos.
\end{enumerate}
\end{warningbox}

\begin{figure}[H]
\centering
\includegraphics[width=0.9\textwidth]{slides-images/slide-16.jpg}
\caption{Retos del reconocimiento de imágenes: cambios de punto de vista, cambios de escala, deformación, oclusión, condiciones de iluminación, fondo desordenado y variación intraclase.\protect\footnotemark}
\end{figure}
\footnotetext{Diapositiva 16.}

\subsection{Aprender jerarquías de características}

La idea central del aprendizaje profundo es: ¿es posible aprender directamente de los datos una jerarquía de características (hierarchy of features), en lugar de diseñarla a mano?

\begin{figure}[H]
\centering
\includegraphics[width=0.85\textwidth]{slides-images/slide-18.jpg}
\caption{Aprender la jerarquía de características a partir de los datos: características de bajo nivel (bordes, puntos oscuros) $\rightarrow$ características de nivel medio (ojos, orejas, nariz) $\rightarrow$ características de alto nivel (estructura facial).\protect\footnotemark}
\end{figure}
\footnotetext{Diapositiva 18. Referencia: Lee y colaboradores, ICML 2009.}

\begin{importantbox}{La idea clave de la jerarquía de características}
Las características visuales se organizan en capas:
\begin{itemize}
    \item \textbf{Características de bajo nivel} (Low-level features): bordes, manchas oscuras y otros elementos visuales básicos
    \item \textbf{Características de nivel medio} (Mid-level features): formadas por la combinación de características de bajo nivel, como ojos y orejas
    \item \textbf{Características de alto nivel} (High-level features): formadas por la combinación de características de nivel medio, como la estructura facial completa
\end{itemize}
El objetivo del aprendizaje profundo es que la red aprenda automáticamente este tipo de representación jerárquica de características.
\end{importantbox}

\subsection{Resumen de la sección}

La extracción manual de características no es escalable debido a los problemas de recursividad y variabilidad. La contribución central del aprendizaje profundo es permitir que la red aprenda automáticamente, a partir de los datos, una representación jerárquica de características, desde la detección de bordes de bajo nivel hasta la comprensión semántica de alto nivel, sin intervención humana en todo el proceso.

\newpage
\section{De las redes totalmente conectadas a la convolución: el aprovechamiento de la estructura espacial}

\subsection{El problema de procesar imágenes con redes totalmente conectadas}

En la primera clase aprendimos las redes neuronales totalmente conectadas (Fully Connected Neural Network). Si se usa directamente una red totalmente conectada para procesar una imagen, hay que \textbf{aplanar} (flatten) primero la imagen bidimensional en un vector unidimensional:

\begin{figure}[H]
\centering
\includegraphics[width=0.9\textwidth]{slides-images/slide-22.jpg}
\caption{El problema de procesar imágenes con redes totalmente conectadas: la operación de aplanado destruye toda la información espacial y produce una cantidad enorme de parámetros.\protect\footnotemark}
\end{figure}
\footnotetext{Diapositivas, página 22.}

\begin{warningbox}{Los dos grandes defectos de las redes totalmente conectadas aplicadas a imágenes}
\begin{enumerate}
    \item \textbf{Pérdida total de la información espacial}: al aplanar la imagen bidimensional en un vector unidimensional, la relación espacial entre los píxeles queda completamente destruida. Píxeles que eran vecinos pueden quedar muy alejados entre sí después del aplanado.
    \item \textbf{Explosión del número de parámetros}: cada píxel se conecta con cada neurona de la capa oculta. Por ejemplo, una imagen en escala de grises de $256 \times 256$ tiene 65,536 píxeles; si la capa oculta tiene 1,024 neuronas, solo la primera capa requiere $65,536 \times 1,024 \approx 67$ millones de parámetros.
\end{enumerate}
\end{warningbox}

La pregunta clave es: \textbf{¿cómo aprovechar la estructura espacial de la entrada para orientar el diseño de la arquitectura de la red?}

\subsection{La idea de la conectividad local}

La idea central de la solución es sencilla: en lugar de que cada neurona se conecte a todos los píxeles de la imagen completa, se conecta únicamente a una \textbf{región local} (local patch) de la imagen de entrada.

\begin{figure}[H]
\centering
\includegraphics[width=0.85\textwidth]{slides-images/slide-23.jpg}
\caption{Aprovechamiento de la estructura espacial: se conecta una región local (patch) de la imagen de entrada a una sola neurona de la capa oculta. Cada neurona solo "ve" una pequeña parte de la imagen.\protect\footnotemark}
\end{figure}
\footnotetext{Diapositivas, página 23.}

Después, mediante una \textbf{ventana deslizante} (sliding window), la misma matriz de pesos se aplica repetidamente en distintas posiciones de la imagen. Esta es la idea central de la operación de \textbf{convolución} (Convolution).

\begin{figure}[H]
\centering
\includegraphics[width=0.85\textwidth]{slides-images/slide-24.jpg}
\caption{Ventana deslizante: se conecta una región local de la entrada a una neurona de la capa oculta, y la ventana deslizante define las conexiones. Pregunta clave: ¿cómo establecer los pesos para detectar una característica específica?\protect\footnotemark}
\end{figure}
\footnotetext{Diapositivas, página 24.}

\begin{importantbox}{Las tres propiedades centrales de la operación de convolución}
\begin{enumerate}
    \item \textbf{Extracción de características locales}: se usa un conjunto de pesos (filtro/filter) para extraer \textbf{características locales}
    \item \textbf{Múltiples filtros}: se usan \textbf{varios filtros} para extraer distintos tipos de características
    \item \textbf{Compartición de parámetros}: los parámetros de cada filtro se \textbf{comparten} espacialmente (spatial sharing), es decir, el mismo filtro se aplica en todas las posiciones de la imagen
\end{enumerate}
\end{importantbox}

\subsection{Resumen de la sección}

Al aplanar la imagen, la red totalmente conectada pierde la información espacial y produce una explosión del número de parámetros. La idea central de la convolución es la conectividad local junto con la compartición de parámetros, lo cual conserva la estructura espacial y reduce drásticamente el número de parámetros.

\newpage
\section{Explicación detallada de la operación de convolución}

\subsection{Estudio de caso: detección de la letra X}

Amini usó un ejemplo didáctico clásico para explicar de manera intuitiva la operación de convolución. Supongamos que queremos determinar si una imagen contiene la letra "X". Aunque dos "X" se vean distintas después de un desplazamiento, un escalamiento o una rotación, la computadora debería poder reconocer ambas como "X".

\begin{figure}[H]
\centering
\includegraphics[width=0.85\textwidth]{slides-images/slide-27.jpg}
\caption{¿X o X? La imagen se representa como una matriz de píxeles, pero la computadora la compara píxel por píxel: incluso para la misma "X", después de un desplazamiento la matriz de valores de píxeles es completamente distinta.\protect\footnotemark}
\end{figure}
\footnotetext{Diapositiva 27. Referencia: Rohrer, "How do CNNs work?".}

Solución: en lugar de comparar a nivel global, se compara a nivel de \textbf{características locales}. Una "X" está formada por varias características: líneas diagonales, puntos de intersección, etcétera. Basta con encontrar estas características locales en ambas imágenes para considerar que ambas son "X".

\begin{figure}[H]
\centering
\includegraphics[width=0.9\textwidth]{slides-images/slide-28.jpg}
\caption{Características de la letra X: segmentos diagonales y un punto de intersección. Aunque las dos X estén ubicadas en posiciones globales distintas, sus características locales coinciden.\protect\footnotemark}
\end{figure}
\footnotetext{Diapositiva 28.}

Estas características locales son los \textbf{filtros} (filters). Cada filtro es una pequeña matriz numérica que representa el patrón que queremos detectar.

\begin{figure}[H]
\centering
\includegraphics[width=0.9\textwidth]{slides-images/slide-29.jpg}
\caption{Tres filtros de 3$\times$3 usados para detectar características de "X": detectan, respectivamente, la diagonal izquierda, la intersección en cruz y la diagonal derecha.\protect\footnotemark}
\end{figure}
\footnotetext{Diapositiva 29.}

\subsection{Definición matemática de la operación de convolución}

Los pasos de la operación de convolución son los siguientes:

\begin{enumerate}
    \item Colocar el filtro en alguna posición de la imagen de entrada
    \item Sobre la región cubierta por el filtro, realizar una \textbf{multiplicación elemento a elemento} (element-wise multiplication)
    \item \textbf{Sumar} todos los productos
    \item Escribir el resultado en la posición correspondiente del \textbf{mapa de características} (feature map)
    \item Deslizar el filtro a la siguiente posición y repetir el proceso anterior
\end{enumerate}

\begin{figure}[H]
\centering
\includegraphics[width=0.9\textwidth]{slides-images/slide-31.jpg}
\caption{Ejemplo de la operación de convolución: convolución de una imagen de entrada de 5$\times$5 con un filtro de 3$\times$3. El filtro se desliza sobre la entrada, se multiplica elemento a elemento y se suma, obteniendo el mapa de características.\protect\footnotemark}
\end{figure}
\footnotetext{Diapositiva 31. Referencia: Karn, "Intuitive CNNs".}

Tomando como ejemplo una imagen de entrada de $5 \times 5$ y un filtro de $3 \times 3$, el proceso de convolución genera un mapa de características de $3 \times 3$:

\begin{figure}[H]
\centering
\includegraphics[width=0.9\textwidth]{slides-images/slide-32.jpg}
\caption{Primer paso de la convolución: el filtro cubre la región de 3$\times$3 de la esquina superior izquierda; al multiplicar elemento a elemento y sumar se obtiene 4, que se escribe en la primera posición del mapa de características.\protect\footnotemark}
\end{figure}
\footnotetext{Diapositiva 32.}

Para el caso general, la operación de convolución puede expresarse con la siguiente fórmula:

$$y_{p,q} = \sum_{i=1}^{k} \sum_{j=1}^{k} w_{i,j} \cdot x_{i+p, j+q} + b$$

donde:
\begin{itemize}
    \item $x$ es la imagen de entrada
    \item $w$ es la matriz de pesos del filtro, de tamaño $k \times k$
    \item $b$ es el término de sesgo
    \item $y_{p,q}$ es el valor de salida del mapa de características en la posición $(p,q)$
    \item $(p,q)$ es la posición de la neurona en la capa oculta
\end{itemize}

\begin{figure}[H]
\centering
\includegraphics[width=0.85\textwidth]{slides-images/slide-40.jpg}
\caption{Resultado completo de la convolución: la convolución de una entrada de 5$\times$5 con un filtro de 3$\times$3 (stride=1) produce un mapa de características de 3$\times$3.\protect\footnotemark}
\end{figure}
\footnotetext{Diapositiva 40.}

\begin{knowledgebox}{La esencia de la convolución es la "suma ponderada"}
Si entendiste la propagación hacia adelante de la capa totalmente conectada de la Clase 1 (realizar una suma ponderada de la entrada y luego sumar un sesgo), entonces lo que hace la capa convolucional es, en esencia, exactamente lo mismo. La única diferencia es que cada neurona de la capa totalmente conectada ve \textbf{toda} la entrada, mientras que cada neurona de la capa convolucional solo ve la entrada dentro de una \textbf{ventana local}, y los pesos de esa ventana se comparten espacialmente.
\end{knowledgebox}

\subsection{Los filtros producen distintos mapas de características}

Filtros con distintos pesos producen mapas de características completamente diferentes. Al cambiar los valores del filtro, se puede lograr:

\begin{figure}[H]
\centering
\includegraphics[width=0.9\textwidth]{slides-images/slide-41.jpg}
\caption{Distintos filtros producen distintos efectos: nitidez (Sharpen), detección de bordes (Edge Detect) y detección de bordes fuerte (Strong Edge Detect).\protect\footnotemark}
\end{figure}
\footnotetext{Diapositiva 41.}

\begin{itemize}
    \item \textbf{Filtro de nitidez}: aumenta el peso del píxel central y disminuye el de los píxeles circundantes, resaltando los detalles de la imagen
    \item \textbf{Filtro de detección de bordes}: es, en esencia, un operador de derivada que detecta los cambios abruptos de brillo de alto a bajo
    \item \textbf{Filtro de detección de bordes fuerte}: un operador de derivada más agresivo, que produce bordes más marcados
\end{itemize}

\begin{knowledgebox}{Del diseño manual al aprendizaje automático de los filtros}
Antes del aprendizaje profundo, durante mucho tiempo se usaron filtros diseñados manualmente (como el operador de Sobel o el operador Laplaciano) para detectar patrones específicos. El avance clave del aprendizaje profundo fue: \textbf{hacer que los pesos del filtro sean parámetros que se pueden aprender}, optimizándolos automáticamente a partir de los datos mediante retropropagación y descenso de gradiente. Así, la red puede descubrir por sí misma las características más útiles para la tarea en cuestión, sin necesidad de diseño humano.
\end{knowledgebox}

\subsection{Cálculo del tamaño de la salida}

El tamaño de salida de la operación de convolución depende de tres hiperparámetros:

\begin{table}[H]
\centering
\begin{tabular}{lp{10cm}}
\toprule
\textbf{Hiperparámetro} & \textbf{Descripción} \\
\midrule
Tamaño del filtro (kernel size) & Las dimensiones espaciales del filtro, por ejemplo $3 \times 3$, $5 \times 5$ \\
Paso (stride) & El número de píxeles que se desplaza el filtro en cada paso. stride=1 significa que se mueve un píxel cada vez; stride=2 significa que salta dos píxeles cada vez \\
Relleno (padding) & Si se agregan ceros en los bordes de la entrada. Si no se rellena, el tamaño de salida se reduce \\
\bottomrule
\end{tabular}
\caption{Hiperparámetros clave de la capa convolucional}
\end{table}

Para una entrada de tamaño $n$, un filtro de tamaño $k$ y un paso de $s$, el tamaño de salida es:

$$\text{output size} = \left\lfloor \frac{n - k}{s} \right\rfloor + 1$$

\subsection{Resumen de la sección}

La operación de convolución desliza un filtro sobre la imagen, realiza multiplicación elemento a elemento y suma, generando un mapa de características. Distintos filtros detectan distintos patrones. En el aprendizaje profundo, los pesos de los filtros se optimizan automáticamente mediante el entrenamiento, en lugar de diseñarse manualmente. Las tres ventajas centrales de la convolución son: conectividad local, uso compartido de parámetros y preservación de la estructura espacial.

\newpage
\section{Arquitecturas de redes neuronales convolucionales (CNN)}

\subsection{Los tres componentes centrales de una CNN}

Habiendo entendido la operación de convolución, podemos construir una red neuronal convolucional completa. Una CNN típica está compuesta por tres tipos de capas que se alternan:

\begin{figure}[H]
\centering
\includegraphics[width=0.9\textwidth]{slides-images/slide-44.jpg}
\caption{Flujo completo de una red CNN de clasificación: imagen de entrada $\rightarrow$ capa convolucional (genera mapas de características) $\rightarrow$ max pooling $\rightarrow$ $\cdots$ $\rightarrow$ capa totalmente conectada $\rightarrow$ salida de clasificación.\protect\footnotemark}
\end{figure}
\footnotetext{Diapositivas, página 44.}

\begin{importantbox}{Las tres operaciones centrales de una CNN}
\begin{enumerate}
    \item \textbf{Convolución} (Convolution): aplica un filtro para generar un mapa de características, extrayendo características locales
    \item \textbf{Activación no lineal} (Non-linearity): normalmente se usa ReLU, que introduce capacidad de representación no lineal
    \item \textbf{Pooling} (Pooling): submuestrea cada mapa de características, reduciendo la dimensión y obteniendo invariancia espacial
\end{enumerate}
Mediante el entrenamiento, la CNN aprende automáticamente los pesos de los filtros de las capas convolucionales.
\end{importantbox}

\subsection{La capa convolucional: conectividad local}

Cada neurona de una capa convolucional se conecta únicamente a una región local de la entrada. Para la neurona en la posición $(p,q)$ de la capa oculta:

\begin{figure}[H]
\centering
\includegraphics[width=0.85\textwidth]{slides-images/slide-46.jpg}
\caption{Conectividad local de la capa convolucional: cada neurona toma su entrada de una región local, calcula la suma ponderada y agrega el sesgo. Su expresión matemática es la misma que la de una capa totalmente conectada, pero su alcance es local.\protect\footnotemark}
\end{figure}
\footnotetext{Diapositivas, página 46.}

$$y_{p,q} = \sum_{i=1}^{k} \sum_{j=1}^{k} w_{ij} \cdot x_{i+p, j+q} + b$$

Este proceso consta de tres pasos: (1) aplicar la ventana de pesos, (2) calcular la combinación lineal, (3) activar mediante una función no lineal.

\subsection{Disposición espacial del volumen de salida}

La salida de cada capa no es un plano bidimensional, sino un \textbf{volumen} (volume) tridimensional, de dimensiones $h \times w \times d$:

\begin{figure}[H]
\centering
\includegraphics[width=0.9\textwidth]{slides-images/slide-47.jpg}
\caption{Dimensiones del volumen de salida de una CNN: $h \times w \times d$, donde $h$ y $w$ son las dimensiones espaciales y $d$ (la profundidad) es igual al número de filtros de esa capa.\protect\footnotemark}
\end{figure}
\footnotetext{Diapositivas, página 47.}

\begin{itemize}
    \item $h$ y $w$: las dimensiones espaciales de la salida (determinadas por el tamaño de la entrada, el tamaño del filtro y el stride)
    \item $d$: la profundidad, igual al \textbf{número de filtros} usado en esa capa
    \item \textbf{Stride} (paso): el tamaño del desplazamiento del filtro al deslizarse
    \item \textbf{Receptive Field} (campo receptivo): la región de la imagen de entrada conectada a un nodo de salida determinado
\end{itemize}

En TensorFlow: \texttt{tf.keras.layers.Conv2D(filters=d, kernel\_size=(h,w), strides=s)}

En PyTorch: \texttt{torch.nn.Conv2d(in\_channels, out\_channels=d, kernel\_size, stride=s)}

\begin{warningbox}{PyTorch requiere especificar el número de canales de entrada}
El \texttt{Conv2D} de TensorFlow/Keras infiere automáticamente la dimensión de entrada, pero el \texttt{Conv2d} de PyTorch exige especificar explícitamente \texttt{in\_channels}. Quienes empiezan suelen olvidar esto y obtienen un error.
\end{warningbox}

\subsection{Función de activación no lineal: ReLU}

Después de cada operación de convolución es necesario aplicar una función de activación no lineal. En las CNN, la más usada es \textbf{ReLU} (Rectified Linear Unit):

\begin{figure}[H]
\centering
\includegraphics[width=0.85\textwidth]{slides-images/slide-48.jpg}
\caption{Función de activación ReLU: reemplaza todos los valores negativos por cero y conserva los valores positivos sin cambio. Es una operación no lineal aplicada píxel por píxel.\protect\footnotemark}
\end{figure}
\footnotetext{Diapositivas, página 48.}

$$g(z) = \max(0, z)$$

La función de ReLU consiste en reemplazar por cero todos los valores negativos del mapa de características, conservando solo los valores positivos (es decir, las regiones donde se detectó la característica). ¿Por qué ReLU resulta especialmente adecuada para tareas con imágenes? Porque cuando un filtro detecta un patrón determinado, la salida es positiva (coincidencia), y cuando no lo detecta, la salida es negativa o cero (no hay coincidencia); ReLU conserva precisamente la señal de «se detectó».

\subsection{Pooling: reducción de dimensión e invariancia espacial}

El \textbf{pooling} (Pooling) es una operación de submuestreo que sigue a la capa convolucional. La más usada es el \textbf{max pooling} (Max Pooling):

\begin{figure}[H]
\centering
\includegraphics[width=0.85\textwidth]{slides-images/slide-49.jpg}
\caption{Ejemplo de max pooling: usando una ventana de 2$\times$2 y un stride de 2, se toma el valor máximo en cada ventana. Una entrada de 4$\times$4 se convierte en una salida de 2$\times$2.\protect\footnotemark}
\end{figure}
\footnotetext{Diapositivas, página 49.}

\begin{importantbox}{Las dos funciones principales del pooling}
\begin{enumerate}
    \item \textbf{Reducción de dimensión} (Reduced dimensionality): reduce el tamaño espacial del mapa de características, disminuyendo así el costo de cómputo y el número de parámetros de las capas siguientes
    \item \textbf{Invariancia espacial} (Spatial invariance): al tomar el valor máximo, la posición exacta de la característica importa menos, lo que aumenta la robustez frente a desplazamientos
\end{enumerate}
\end{importantbox}

Amini menciona que, además del max pooling, se puede plantear una pregunta más general: «¿de qué otras formas se puede lograr el submuestreo y a la vez mantener la invariancia espacial?». Esto sugiere otras técnicas que quizá se presenten en clases posteriores, como la convolución con stride (strided convolution).

\subsection{Aprendizaje de representaciones en CNN profundas}

Al apilar varios módulos de «convolución + ReLU + pooling», la CNN es capaz de aprender representaciones de características cada vez más abstractas:

\begin{figure}[H]
\centering
\includegraphics[width=0.85\textwidth]{slides-images/slide-50.jpg}
\caption{Aprendizaje de representaciones en una CNN profunda: la primera capa aprende características de bajo nivel como bordes y manchas oscuras, la segunda las combina en características de nivel medio como ojos y nariz, y la tercera las combina en características de alto nivel como la estructura facial completa.\protect\footnotemark}
\end{figure}
\footnotetext{Diapositivas, página 50. Referencia: Lee y otros, ICML 2009.}

\begin{knowledgebox}{Composición jerárquica de características}
Las CNN profundas son poderosas porque cada capa realiza una \textbf{composición} sobre las características detectadas por la capa anterior:
\begin{itemize}
    \item La primera capa convolucional detecta características primitivas (como bordes en distintas direcciones)
    \item Después del pooling, el mapa de características se reduce, lo que equivale a observar a una escala espacial mayor
    \item La segunda capa convolucional detecta sobre el mapa de características reducido, lo que equivale a combinar varias características de bajo nivel en características de nivel medio
    \item Al repetirse este proceso, las características se vuelven cada vez más abstractas y semánticamente más ricas
\end{itemize}
Este aprendizaje jerárquico de características, de lo simple a lo complejo, es precisamente lo que significa la palabra «profundo» en aprendizaje profundo.
\end{knowledgebox}

\subsection{Resumen de la sección}

Una red CNN de clasificación completa está compuesta por una parte de aprendizaje de características y una parte de clasificación. La parte de aprendizaje de características, alternando capas convolucionales (que extraen características locales), activación no lineal (que aporta capacidad de representación) y capas de pooling (que reducen la dimensión y aumentan la invariancia), construye capa por capa representaciones de características cada vez más abstractas. La parte de clasificación aplana las características aprendidas y las pasa por una capa totalmente conectada para obtener las probabilidades de clase.

\newpage
\section{Flujo completo de la CNN: de la extracción de características a la clasificación}

\subsection{Panorama general de la arquitectura completa}

Una red CNN de clasificación completa puede dividirse en dos etapas:

\begin{figure}[H]
\centering
\includegraphics[width=0.9\textwidth]{slides-images/slide-51.jpg}
\caption{Flujo completo de la CNN: la primera mitad es el \textbf{aprendizaje de características} (repetición de convolución + ReLU + agrupamiento), la segunda mitad es la \textbf{clasificación} (aplanado + capa totalmente conectada + Softmax).\protect\footnotemark}
\end{figure}
\footnotetext{Diapositiva 51.}

\begin{enumerate}
    \item \textbf{Etapa de aprendizaje de características}:
    \begin{itemize}
        \item Aprender las características de la imagen de entrada mediante \textbf{convolución}
        \item Introducir no linealidad mediante \textbf{activación no lineal} (porque los datos del mundo real son no lineales)
        \item Reducir la dimensionalidad y preservar la invarianza espacial mediante \textbf{agrupamiento}
        \item Repetir el proceso anterior varias veces
    \end{itemize}
    \item \textbf{Etapa de clasificación}:
    \begin{itemize}
        \item \textbf{Aplanar} (flatten) el mapa de características final en un vector unidimensional
        \item Clasificar mediante una \textbf{capa totalmente conectada}
        \item Usar \textbf{Softmax} para obtener la probabilidad de cada clase
    \end{itemize}
\end{enumerate}

\subsection{Cabeza de clasificación: salida Softmax}

Después de que las capas de convolución y agrupamiento de la CNN producen una representación de características de alto nivel, esta se proyecta a las distintas clases mediante una capa totalmente conectada, y finalmente la función Softmax produce una distribución de probabilidad:

\begin{figure}[H]
\centering
\includegraphics[width=0.9\textwidth]{slides-images/slide-52.jpg}
\caption{Cabeza de clasificación de la CNN: las capas CONV y POOL producen características de alto nivel que, tras el aplanado, pasan por una capa totalmente conectada, y finalmente Softmax produce la probabilidad de cada clase.\protect\footnotemark}
\end{figure}
\footnotetext{Diapositiva 52.}

$$\text{softmax}(y_i) = \frac{e^{y_i}}{\sum_j e^{y_j}}$$

donde $y_i$ es la puntuación cruda (logit) de la $i$-ésima clase producida por la capa totalmente conectada; Softmax la convierte en un valor de probabilidad, y la suma de las probabilidades de todas las clases es igual a 1.

\subsection{Implementación en código: TensorFlow y PyTorch}

\begin{figure}[H]
\centering
\includegraphics[width=0.9\textwidth]{slides-images/slide-53.jpg}
\caption{Implementación de una CNN en TensorFlow/Keras: dos capas de convolución (32 y 64 filtros) + aplanado + capa totalmente conectada + salida Softmax de 10 clases.\protect\footnotemark}
\end{figure}
\footnotetext{Diapositiva 53.}

\begin{lstlisting}[caption={Implementación de CNN en TensorFlow/Keras}]
import tensorflow as tf

def generate_model():
    model = tf.keras.Sequential([
        # First convolutional layer
        tf.keras.layers.Conv2D(32, filter_size=3, activation='relu'),
        tf.keras.layers.MaxPool2D(pool_size=2, strides=2),

        # Second convolutional layer
        tf.keras.layers.Conv2D(64, filter_size=3, activation='relu'),
        tf.keras.layers.MaxPool2D(pool_size=2, strides=2),

        # Fully connected classifier
        tf.keras.layers.Flatten(),
        tf.keras.layers.Dense(1024, activation='relu'),
        tf.keras.layers.Dense(10, activation='softmax')  # 10 outputs
    ])
    return model
\end{lstlisting}

\begin{lstlisting}[caption={Implementación de CNN en PyTorch}]
import torch
import torch.nn as nn

def generate_model():
    model = nn.Sequential(
        # First and second convolutional layer
        nn.Conv2d(in_channels=3, out_channels=32, kernel_size=3),
        nn.ReLU(),
        nn.MaxPool2d(kernel_size=2, stride=2),

        nn.Conv2d(in_channels=32, out_channels=64, kernel_size=3),
        nn.ReLU(),
        nn.MaxPool2d(kernel_size=2, stride=2),

        # Fully connected classifier
        nn.Flatten(),
        nn.Linear(64*6*6, 1024),  # flattened dim after 2 conv layers
        nn.ReLU(),
        nn.Linear(1024, 10),      # 10 outputs
    )
    return model
\end{lstlisting}

\begin{knowledgebox}{El "arte y la ciencia" de elegir los hiperparámetros}
Amini reconoce abiertamente que la elección de los hiperparámetros de la CNN (número y tamaño de los filtros, número de capas, etc.) es "más un arte que una ciencia" (more art than science), pero también existen algunas pautas empíricas:
\begin{itemize}
    \item El tamaño de los filtros suele ser pequeño ($3 \times 3$ o $5 \times 5$), porque las características locales no requieren un campo receptivo grande
    \item El número de filtros suele \textbf{aumentar capa por capa} (por ejemplo, $32 \rightarrow 64 \rightarrow 128$), porque las capas más altas necesitan representar más tipos de características compuestas
    \item Las dimensiones espaciales \textbf{disminuyen capa por capa} mediante el agrupamiento, en complemento con el aumento del número de filtros
\end{itemize}
\end{knowledgebox}

\subsection{Resumen de la sección}

Una red CNN de clasificación completa se divide en dos etapas: aprendizaje de características y clasificación. El aprendizaje de características extrae características abstractas capa por capa mediante el apilamiento de convolución + ReLU + agrupamiento; la etapa de clasificación proyecta las características a probabilidades de clase mediante el aplanado y una capa totalmente conectada. Tanto TensorFlow como PyTorch ofrecen API sencillas para construir una CNN.

\newpage
\section{Aplicaciones extensas de las CNN}

\subsection{Más allá de la clasificación: una arquitectura de propósito general}

La parte de aprendizaje de características de las CNN es altamente general: la representación de características aprendida puede usarse para múltiples tareas posteriores distintas:

\begin{figure}[H]
\centering
\includegraphics[width=0.85\textwidth]{slides-images/slide-56.jpg}
\caption{Las CNN como arquitectura de propósito general: el mismo módulo de aprendizaje de características puede sustentar clasificación, detección de objetos, segmentación semántica, control probabilístico y muchas otras tareas.\protect\footnotemark}
\end{figure}
\footnotetext{Diapositiva 56.}

\begin{itemize}
    \item \textbf{Clasificación} (Classification): determinar a qué categoría pertenece toda la imagen
    \item \textbf{Detección de objetos} (Object Detection): localizar y clasificar todos los objetos de la imagen
    \item \textbf{Segmentación semántica} (Segmentation): clasificar cada píxel
    \item \textbf{Control probabilístico} (Probabilistic Control): por ejemplo, en conducción autónoma, ir de la imagen al ángulo del volante
\end{itemize}

\subsection{Aplicación uno: detección de cáncer de mama en imágenes médicas}

Las CNN han logrado resultados revolucionarios en el análisis de imágenes médicas. Amini presentó un estudio publicado en Nature que demuestra que un sistema de IA basado en CNN \textbf{superó el desempeño de radiólogos especialistas} en la detección de cáncer de mama:

\begin{figure}[H]
\centering
\includegraphics[width=0.9\textwidth]{slides-images/slide-57.jpg}
\caption{CNN aplicadas a la detección de cáncer de mama: la curva ROC muestra que la sensibilidad y la especificidad del sistema de IA superan a las de los radiólogos. La imagen de la derecha muestra un caso de lesión mamaria detectada por la IA que un médico humano había pasado por alto.\protect\footnotemark}
\end{figure}
\footnotetext{Diapositiva 57. Referencia: McKinney y colaboradores, Nature 2020.}

\begin{importantbox}{El valor de las CNN en las imágenes médicas}
\begin{itemize}
    \item En la evaluación de la detección de cáncer de mama a 2 años y a 1 año en Estados Unidos, la curva ROC del sistema de IA superó a la del grupo de radiólogos
    \item La IA puede detectar lesiones diminutas que los médicos humanos pasan por alto, lo cual tiene un valor clínico auxiliar importante
    \item Métodos de CNN similares ya se han extendido a la detección de cáncer de piel (Esteva y colaboradores, Nature 2017) y al diagnóstico por imágenes de COVID-19
\end{itemize}
\end{importantbox}

\subsection{Aplicación dos: detección de objetos}

La \textbf{detección de objetos} (Object Detection) va más allá de la clasificación: no solo hay que identificar la categoría del objeto, sino también dar su ubicación precisa (bounding box).

\begin{figure}[H]
\centering
\includegraphics[width=0.85\textwidth]{slides-images/slide-58.jpg}
\caption{Clasificación contra detección de objetos: la clasificación solo produce una etiqueta de categoría (por ejemplo, Taxi); la detección de objetos produce a la vez la categoría y la posición del cuadro delimitador $(x, y, w, h)$.\protect\footnotemark}
\end{figure}
\footnotetext{Diapositiva 58.}

El formato de salida de la detección de objetos es un conjunto $(class, x, y, w, h)$ por cada objeto detectado:
\begin{itemize}
    \item $class$: categoría del objeto
    \item $(x, y)$: coordenadas de la esquina superior izquierda del cuadro delimitador
    \item $(w, h)$: ancho y alto del cuadro delimitador
\end{itemize}

\begin{figure}[H]
\centering
\includegraphics[width=0.85\textwidth]{slides-images/slide-59.jpg}
\caption{Detección de múltiples objetos: el modelo debe producir la categoría y el cuadro delimitador de \textbf{cada} objeto de la escena, y el número de objetos no es fijo.\protect\footnotemark}
\end{figure}
\footnotetext{Diapositiva 59.}

\begin{warningbox}{Dificultades de la detección de objetos}
La detección de objetos es mucho más difícil que la clasificación, sobre todo por lo siguiente:
\begin{enumerate}
    \item El cuadro delimitador puede aparecer en \textbf{cualquier posición} de la imagen
    \item El cuadro delimitador puede ser de \textbf{cualquier tamaño}
    \item El \textbf{número de objetos} de la escena no es fijo: puede ser cero, uno o varios
    \item Distintos objetos pueden tener distintas categorías
\end{enumerate}
Esto hace que la dimensión del espacio de salida ya no sea fija, por lo que se necesita un diseño de red especial para manejarla.
\end{warningbox}

\subsubsection{Método ingenuo: ventana deslizante}

El método más intuitivo consiste en probar en la imagen todos los cuadros delimitadores posibles, en todas las posiciones y tamaños, y clasificar con una CNN la región de imagen dentro de cada cuadro. Pero el costo de cómputo de este método es catastrófico: el número de ventanas posibles crece de manera exponencial con la posición, el tamaño y la proporción de aspecto.

\begin{figure}[H]
\centering
\includegraphics[width=0.85\textwidth]{slides-images/slide-60.jpg}
\caption{Detección de objetos ingenua: se recorren todas las ventanas posibles y cada una se envía a una CNN para clasificarla. Problema: hay demasiadas ventanas, y las combinaciones de distintas posiciones, escalas y tamaños producen una cantidad enorme de regiones candidatas.\protect\footnotemark}
\end{figure}
\footnotetext{Diapositiva 60.}

\subsubsection{La familia R-CNN}

Para resolver el problema de eficiencia, los investigadores propusieron la familia de métodos R-CNN (Region-based CNN). La idea central es generar primero un número reducido de regiones candidatas (region proposals) con métodos como la búsqueda selectiva (Selective Search), y clasificar con una CNN solo esas regiones candidatas. Los desarrollos posteriores incluyen Fast R-CNN y Faster R-CNN, que integraron aún más una red de propuesta de regiones (Region Proposal Network, RPN) dentro de la CNN, para lograr un entrenamiento de extremo a extremo.

\begin{knowledgebox}{La evolución de R-CNN a Faster R-CNN}
\begin{itemize}
    \item \textbf{R-CNN} (2014): la búsqueda selectiva genera cerca de 2,000 regiones candidatas, cada una pasa de manera independiente por una CNN para extraer características, y después se clasifica con una SVM. Es muy lento.
    \item \textbf{Fast R-CNN} (2015): primero se hace una sola extracción de características con CNN sobre toda la imagen, y luego se recortan las regiones candidatas a partir del mapa de características. Al compartir el cálculo de características, la velocidad mejora considerablemente.
    \item \textbf{Faster R-CNN} (2015): introduce una Region Proposal Network (RPN) que genera las regiones candidatas directamente sobre el mapa de características, con entrenamiento completamente de extremo a extremo.
\end{itemize}
\end{knowledgebox}

\subsection{Resumen de la sección}

El módulo de aprendizaje de características de las CNN es general y puede sustentar múltiples tareas de visión, como clasificación, detección de objetos y segmentación. En las imágenes médicas, las CNN ya alcanzan e incluso superan el nivel de los expertos humanos. La detección de objetos requiere predecir a la vez la categoría y la posición; el método ingenuo de ventana deslizante es poco eficiente, y la familia R-CNN mejoró considerablemente la eficiencia mediante el mecanismo de propuesta de regiones.

\newpage
\section{Sesión de laboratorio: Lab de detección de rostros}

Amini presentó en clase el laboratorio complementario de esta clase (Lab 2): la tarea de detección de rostros. Este laboratorio permite a los estudiantes construir con sus propias manos una CNN para detectar rostros en imágenes, una excelente oportunidad para llevar a la práctica la teoría de esta clase.

Los pasos principales del laboratorio incluyen:
\begin{enumerate}
    \item Construir la arquitectura de la red neuronal convolucional (incluyendo capas convolucionales, activación ReLU, capas de pooling y una cabeza de clasificación totalmente conectada)
    \item Entrenar el modelo con un conjunto de datos de imágenes de rostros
    \item Evaluar el desempeño de detección del modelo en imágenes nuevas
    \item Analizar la visualización de las características aprendidas por el modelo
\end{enumerate}

\begin{knowledgebox}{El significado histórico de la detección de rostros}
La detección de rostros es una de las primeras tareas de visión por computadora en obtener una aplicación comercial a gran escala. Del detector de rostros Viola-Jones (2001, basado en características Haar diseñadas manualmente y en un clasificador AdaBoost) hasta los métodos modernos de aprendizaje profundo, la precisión y la velocidad de la detección de rostros han dado un salto cualitativo. Hoy en día, la tecnología de detección de rostros se ha integrado en aplicaciones cotidianas como el desbloqueo de teléfonos, los filtros de redes sociales y la vigilancia de seguridad.
\end{knowledgebox}

\begin{warningbox}{Consideraciones éticas del reconocimiento facial}
Aunque las tecnologías de detección y reconocimiento facial son muy poderosas, también han suscitado graves problemas de privacidad y ética. Estos incluyen: sesgos raciales y de género (el modelo se desempeña peor en ciertos grupos), abuso de la vigilancia y falsificaciones mediante deepfake, entre otros. Al desarrollar y desplegar este tipo de tecnología, es indispensable considerar seriamente su impacto social.
\end{warningbox}

\newpage
\section{Síntesis y ampliación}

\subsection{Puntos clave de esta clase}

\begin{enumerate}
    \item \textbf{Una imagen es números}: la computadora representa una imagen mediante una matriz numérica de $H \times W \times C$, donde $C$ es el número de canales (imagen en escala de grises $C=1$, imagen a color $C=3$).

    \item \textbf{La extracción de características es el núcleo del reconocimiento visual}: para reconocer un objeto es necesario extraer características que distingan entre categorías distintas. El diseño manual de características enfrenta el dilema de la recursividad y la variabilidad; el aprendizaje profundo resuelve este problema aprendiendo automáticamente una jerarquía de características a partir de los datos.

    \item \textbf{Las redes totalmente conectadas no son adecuadas para imágenes}: la operación de aplanado hace perder la información espacial y el número de parámetros resulta excesivo.

    \item \textbf{Las tres ventajas principales de la operación de convolución}:
    \begin{itemize}
        \item conectividad local: conserva la estructura espacial
        \item compartición de parámetros: reduce drásticamente el número de parámetros
        \item equivarianza a la traslación: un mismo patrón se detecta sin importar en qué posición de la imagen aparezca
    \end{itemize}

    \item \textbf{Arquitectura de las CNN}: el apilamiento de capas convolucionales, ReLU y capas de agrupamiento (pooling) constituye el módulo de aprendizaje de características, mientras que las capas totalmente conectadas y softmax constituyen el módulo de clasificación.

    \item \textbf{Las CNN son una arquitectura visual de propósito general}: el mismo marco de aprendizaje de características puede emplearse en clasificación, detección de objetos, segmentación semántica y muchas otras tareas.
\end{enumerate}

\subsection{Relación con las dos clases anteriores}

\begin{table}[H]
\centering
\begin{tabular}{p{3cm}p{5cm}p{5cm}}
\toprule
\textbf{Aspecto} & \textbf{Lecture 1 (red totalmente conectada)} & \textbf{Lecture 3 (CNN)} \\
\midrule
Forma de la entrada & vector de características unidimensional & tensor de imagen bidimensional o tridimensional \\
Forma de conexión & totalmente conectada (todas las entradas a todas las neuronas) & conectividad local (cada neurona solo observa una región local) \\
Compartición de parámetros & ninguna & el filtro se comparte a lo largo del espacio \\
Información espacial & no se conserva & se conserva y se aprovecha \\
Operación central & multiplicación de matrices & operación de convolución \\
\bottomrule
\end{tabular}
\caption{Comparación entre la red totalmente conectada y la CNN}
\end{table}

\begin{importantbox}{Filosofía central: dejar que la estructura de los datos guíe el diseño de la arquitectura}
La filosofía de diseño de las CNN consiste en \textbf{aprovechar la estructura interna de los datos para diseñar la arquitectura de la red}. Una imagen tiene localidad espacial: la relación entre píxeles vecinos es más estrecha que entre píxeles distantes; una imagen tiene invariancia a la traslación: un mismo patrón puede aparecer en cualquier posición. La operación de convolución aprovecha perfectamente estos dos conocimientos previos. Esta filosofía se repite una y otra vez en el aprendizaje profundo: las RNN aprovechan la estructura temporal de las secuencias, y el mecanismo de atención de los Transformer aprovecha la estructura de relaciones entre elementos.
\end{importantbox}

\subsection{Lecturas adicionales}

\begin{itemize}
    \item Sitio oficial del curso MIT 6.S191: \url{http://introtodeeplearning.com}
    \item LeCun et al. (1998), "Gradient-Based Learning Applied to Document Recognition" — obra fundacional de las CNN
    \item Krizhevsky et al. (2012), "ImageNet Classification with Deep Convolutional Neural Networks" (AlexNet) — hito del aprendizaje profundo en el campo de la visión por computadora
    \item Lee et al. (2009), ICML — artículo original de la visualización de la jerarquía de características presentada en esta clase
    \item McKinney et al. (2020), Nature — estudio en el que una CNN supera a médicos humanos en la detección de cáncer de mama
    \item Esteva et al. (2017), Nature — estudio sobre el uso de CNN para el diagnóstico de cáncer de piel
    \item Rohrer, "How do CNNs work?" — referencia del caso de detección de la letra X presentado en esta clase
    \item Karn, "An Intuitive Explanation of CNNs" — referencia de la animación de la operación de convolución presentada en esta clase
\end{itemize}

%% --- Fin del contenido --- %%

\end{document}
